% ========================================================
% BÀI TẬP TOÁN 5 - KẾT NỐI TRI THỨC VỚI CUỘC SỐNG
% Bài 24: Luyện tập chung
% Định dạng: BTTeX Part III - Trắc nghiệm điền số ngắn
% ========================================================

% Câu 1
\begin{ex}
  Tính giá trị biểu thức: $12{,}5 \times 4 - 35 : 0{,}7$.
  \shortans{0}
  \loigiai{
    $12{,}5 \times 4 = 50$; $35 : 0{,}7 = 50$. $50 - 50 = 0$.
  }
\end{ex}

% Câu 2
\begin{ex}
  Tính giá trị biểu thức: $(45{,}8 + 14{,}2) : 2{,}5$.
  \shortans{24}
  \loigiai{
    $60 : 2{,}5 = 24$.
  }
\end{ex}

% Câu 3
\begin{ex}
  Tính nhanh: $4{,}5 \times 36 + 4{,}5 \times 64$.
  \shortans{450}
  \loigiai{
    $4{,}5 \times 100 = 450$.
  }
\end{ex}

% Câu 4
\begin{ex}
  Tính nhanh: $12{,}8 \times 75 - 12{,}8 \times 65$.
  \shortans{128}
  \loigiai{
    $12{,}8 \times 10 = 128$.
  }
\end{ex}

% Câu 5
\begin{ex}
  Tìm $x$ biết: $x \times 3{,}5 = 140$.
  \shortans{40}
  \loigiai{
    $x = 140 : 3{,}5 = 40$.
  }
\end{ex}

% Câu 6
\begin{ex}
  Tìm $x$ biết: $x : 1{,}25 = 16$.
  \shortans{20}
  \loigiai{
    $x = 16 \times 1{,}25 = 20$.
  }
\end{ex}

% Câu 7
\begin{ex}
  Một hình chữ nhật có chu vi $30\text{ m}$, chiều dài hơn chiều rộng $3\text{ m}$. Diện tích hình chữ nhật là bao nhiêu mét vuông? (Nhân 4 lên: hoặc số đẹp)
  \shortans{54}
  \loigiai{
    Nửa chu vi là 15. Chiều dài là $(15 + 3) : 2 = 9\text{ m}$. Chiều rộng là 6 m. Diện tích là $9 \times 6 = 54\text{ m}^2$.
  }
\end{ex}

% Câu 8
\begin{ex}
  Mua $5\text{ m}$ vải hết $125\,000$ đồng. Hỏi mua $8\text{ m}$ vải như thế hết bao nhiêu nghìn đồng?
  \shortans{200}
  \loigiai{
    1 m hết $125 : 5 = 25$ nghìn. 8 m hết $25 \times 8 = 200$ nghìn đồng.
  }
\end{ex}

% Câu 9
\begin{ex}
  Một người đi xe máy trong $2{,}5\text{ giờ}$ đi được $100\text{ km}$. Hỏi vận tốc xe máy là bao nhiêu km/h?
  \shortans{40}
  \loigiai{
    $100 : 2{,}5 = 40\text{ km/h}$.
  }
\end{ex}

% Câu 10
\begin{ex}
  Một mảnh đất hình chữ nhật có diện tích $120\text{ m}^2$, chiều rộng $7{,}5\text{ m}$. Chiều dài mảnh đất là bao nhiêu mét?
  \shortans{16}
  \loigiai{
    Chiều dài là: $120 : 7{,}5 = 16\text{ m}$.
  }
\end{ex}

% Câu 11
\begin{ex}
  Tính giá trị biểu thức: $100 - (15{,}5 + 24{,}5) \times 2$.
  \shortans{20}
  \loigiai{
    $100 - 40 \times 2 = 100 - 80 = 20$.
  }
\end{ex}

% Câu 12
\begin{ex}
  Tính nhanh: $2{,}5 \times 8 \times 5$.
  \shortans{100}
  \loigiai{
    $20 \times 5 = 100$.
  }
\end{ex}

% Câu 13
\begin{ex}
  Một thùng dầu có $50$ lít. Người ta lấy ra $3$ lần, mỗi lần $12{,}5$ lít. Trong thùng còn lại bao nhiêu lít dầu? (Nhân đôi số dầu còn lại).
  \shortans{25}
  \loigiai{
    Lấy ra: $12{,}5 \times 3 = 37{,}5$ lít. Còn lại: $50 - 37{,}5 = 12{,}5$ lít. Nhân đôi: 25.
  }
\end{ex}

% Câu 14
\begin{ex}
  Một xe ô tô chở $25$ bao gạo, mỗi bao cân nặng $40{,}4\text{ kg}$. Tổng khối lượng gạo xe chở là bao nhiêu ki-lô-gam?
  \shortans{1010}
  \loigiai{
    $25 \times 40{,}4 = 1\,010\text{ kg}$.
  }
\end{ex}

% Câu 15
\begin{ex}
  Một công nhân trong $4$ ngày dệt được $72\text{ m}$ vải. Hỏi trong $6{,}5$ ngày người đó dệt được bao nhiêu mét vải?
  \shortans{117}
  \loigiai{
    1 ngày dệt $72 : 4 = 18\text{ m}$. Trong $6{,}5$ ngày dệt: $18 \times 6{,}5 = 117\text{ m}$.
  }
\end{ex}

% Câu 16
\begin{ex}
  Tính giá trị: $(2{,}4 \times 5) + (3{,}6 : 0{,}2)$.
  \shortans{30}
  \loigiai{
    $12 + 18 = 30$.
  }
\end{ex}

% Câu 17
\begin{ex}
  Tìm số tự nhiên $a$ biết: $a : 0{,}5 + a : 0{,}2 = 70$.
  \shortans{10}
  \loigiai{
    $2a + 5a = 70 \implies 7a = 70 \implies a = 10$.
  }
\end{ex}

% Câu 18
\begin{ex}
  Một can dầu hỏa nặng $12\text{ kg}$. Sau khi dùng hết nửa lượng dầu thì can còn nặng $6{,}5\text{ kg}$. Khối lượng vỏ can là bao nhiêu ki-lô-gam?
  \shortans{1}
  \loigiai{
    Nửa lượng dầu nặng: $12 - 6{,}5 = 5{,}5\text{ kg}$. Cả lượng dầu nặng: $5{,}5 \times 2 = 11\text{ kg}$. Vỏ can nặng: $12 - 11 = 1\text{ kg}$.
  }
\end{ex}

% Câu 19
\begin{ex}
  Tính: $1{,}25 \times 4 \times 1{,}25 \times 4$.
  \shortans{25}
  \loigiai{
    $5 \times 5 = 25$.
  }
\end{ex}

% Câu 20
\begin{ex}
  Một hình vuông có cạnh $3{,}5\text{ m}$. Chu vi hình vuông đó là bao nhiêu mét?
  \shortans{14}
  \loigiai{
    Chu vi là $3{,}5 \times 4 = 14\text{ m}$.
  }
\end{ex}

% Câu 21
\begin{ex}
  Một tam giác có diện tích $45\text{ cm}^2$, đáy dài $10\text{ cm}$. Chiều cao tương ứng là bao nhiêu xăng-ti-mét?
  \shortans{9}
  \loigiai{
    Chiều cao là $(45 \times 2) : 10 = 9\text{ cm}$.
  }
\end{ex}

% Câu 22
\begin{ex}
  Tính nhanh: $0{,}4 \times 370 + 0{,}4 \times 630$.
  \shortans{400}
  \loigiai{
    $0{,}4 \times 1\,000 = 400$.
  }
\end{ex}

% Câu 23
\begin{ex}
  Một cửa hàng bán được $45$ chai nước giặt loại $1{,}8$ lít. Tổng số lít nước giặt đã bán là bao nhiêu lít?
  \shortans{81}
  \loigiai{
    $45 \times 1{,}8 = 81$ lít.
  }
\end{ex}

% Câu 24
\begin{ex}
  Tìm số tự nhiên $x$ biết $15{,}4 < x \times 3 < 20{,}5$.
  \shortans{6}
  \loigiai{
    $5{,}13 < x < 6{,}83$. Số tự nhiên duy nhất là 6.
  }
\end{ex}

% Câu 25
\begin{ex}
  Khi nhân một số với $2{,}5$, một bạn đã nhân nhầm với $25$ nên được kết quả là $150$. Kết quả đúng của phép tính là bao nhiêu?
  \shortans{15}
  \loigiai{
    Số đó là $150 : 25 = 6$. Kết quả đúng là $6 \times 2{,}5 = 15$.
  }
\end{ex}
